\begin{table}[htbp]
  \centering
  \small
  \caption{Estimands defined in this study, consolidating the identification qualifications
    stated across Sections~\ref{sec:methodology} and~\ref{sec:results}. Each quantity is
    \emph{descriptive} of the treatment as defined; only relative within-condition differences are
    meant to travel, and rankings are configuration- and version-specific.}
  \label{tab:estimands}
  \begin{adjustbox}{max width=\textwidth}
  \begin{tabular}{@{}p{2.7cm} p{4.0cm} p{4.0cm} p{2.7cm}@{}}
    \toprule
    \textbf{Estimand / quantity} & \textbf{What it identifies} & \textbf{What it does \emph{not} identify} & \textbf{Where reported} \\
    \midrule
    Documentation-selected implementation-and-strategy effect (RQ1, primary) &
    Performance a developer following each library's documentation obtains (throughput, p99), as defined &
    Intrinsic library overhead; the most common or most performant usage &
    Tables~\ref{tab:deep_fetch},~\ref{tab:write} \\
    \addlinespace
    Same-SQL (raw-path) standardized contrast (diagnostic) &
    Residual spread once every layer runs common SQL through its raw facility &
    A causal decomposition or a bound on the library effect; the share from SQL formulation or query count &
    Supplement Table~S42 \\
    \addlinespace
    Performance-conscious deep-fetch regime &
    Speed recoverable within a narrow byte-equivalent tuning budget (faster documented loading strategy) &
    Gains from SQL rewrites, caching, or schema change; non-drop-in strategies (excluded) &
    Table~\ref{tab:deepfetch_regimes} \\
    \addlinespace
    Per-pattern capacity (saturating throughput) &
    Each pattern's throughput knee per layer and engine, locating relative utilization at the operating point &
    Tail latency; a throughput quantity, not a demand- or utilization-dependent comparison on its own &
    Supplement Table~S35, Figure~S2 \\
    \addlinespace
    High-load p99 under equal closed-loop demand (primary tail) &
    The tail a deployment sees at fixed 50-connection demand, including acquisition queueing &
    The tail at comparable utilization; the pooled cross-run p99 &
    Table~\ref{tab:deep_fetch} \\
    \addlinespace
    Matched-utilization tail (p99 at matched fractions of own capacity) &
    p99 when each layer is offered equal fractions (50/70/85/95\%) of its \emph{own} saturating throughput &
    The tail under equal external demand; the capacity-and-queueing component it removes &
    Supplement Table~S43 \\
    \addlinespace
    Layer$\times$engine interaction magnitude (RQ2) &
    The spread of each layer's PostgreSQL$\div$MySQL advantage and the concrete rank reversals &
    The interaction's size from the permutation test (bounds existence only); its sign (every layer faster on PostgreSQL) &
    Supplement Table~S28; \S\ref{sec:results} (RQ2) \\
    \addlinespace
    Cross-engine rank correlations ($n=7$, coarse descriptive) &
    Descriptive agreement of the layer ordering across engines (Spearman~$\rho$, Kendall~$\tau$) &
    Reliable cross-engine transfer at $n=7$; the RQ2 finding, which rests on the reversals and interaction spread &
    \S\ref{sec:results} (RQ2) \\
    \bottomrule
  \end{tabular}
  \end{adjustbox}
\end{table}
